%% kapa: exam=1 round=35 year=2024 period=2 form=A subject=appraisal_law subject_ko=감정평가관계법규 seq=26 number=26
%% source: Q-Net 감정평가사 기출문제 (cst00309) artlSeq=5213229 file=raw/1008/5213229/2교시 문제지 원본.pdf
%% conversion: kordoc 4.24.51 -> tools/gen1.py (auto); review=proofread-sonnet; flags=none
\begin{question}{26}
건축법령상 특별가로구역에 관한 조문의 내용이다. ( )에 들어갈 내용으로 옳은 것은?
\begin{framedtext}
국토교통부장관 및 허가권자는 「건축법」 및 관계 법령에 따라 일부 규정을 적용하지 아니하거나 완화하여 적용할 수 있도록 ( ㄱ )에서 ( ㄴ )에 접한 대지의 일정 구역을 특별가로구역으로 지정할 수 있다.
\end{framedtext}
\begin{choices}
  \item ㄱ: 개발진흥지구, ㄴ: 허가권자가 리모델링 활성화가 필요하다고 인정하여 지정ㆍ공고한 지역 안의 도로
  \item ㄱ: 경관지구, ㄴ: 「지역문화진흥법」에 따른 문화지구 안의 도로
  \item ㄱ: 개발진흥지구, ㄴ: 보행자전용도로로서 도시미관 개선을 위하여 허가권자가 건축조례로 정하는 도로
  \item ㄱ: 경관지구, ㄴ: 「도시 및 주거환경정비법」에 따른 정비구역 안의 도로
  \item ㄱ: 개발진흥지구, ㄴ: 건축선을 후퇴한 대지에 접한 도로로서 허가권자가 건축조례로 정하는 도로
\end{choices}
\end{question}
\answer{2}
